\akapit Przeprowadzone symulacje numeryczne pozwalają na wyciągnięcie ogólnych wniosków dotyczących wrażliwości systemu sterowania bezzałogowego statku powietrznego na różnego rodzaju zakłócenia i ograniczenia sprzętowe. Analiza lotu swobodnego potwierdziła poprawność modelu oraz stabilność statyczną i dynamiczną płatowca w warunkach nominalnych. Wprowadzenie ograniczeń w układzie wykonawczym – takich jak skończona prędkość serwomechanizmów – uwidoczniło istotną zwłokę fazową w torze sterowania, co bezpośrednio skutkuje gorszym tłumieniem odchyłek i wzrostem błędu śledzenia trajektorii. Z kolei obecność luzów kinematycznych generuje strefę martwą, która uniemożliwia precyzyjną korektę położenia, prowadząc do charakterystycznego oscylacyjnego „pływania” modelu wokół zadanej ścieżki.

\akapit Szczególnie krytycznym czynnikiem dla stabilności całego systemu okazały się opóźnienia transportowe w pętli sprzężenia zwrotnego. Nawet umiarkowane opóźnienie drastycznie obniża zapas stabilności regulatora, co prowadzi do gwałtownych, narastających oscylacji i uniemożliwia prawidłowe wykonanie misji. Z drugiej strony, dodanie szumów pomiarowych, choć wywołuje widoczną nerwowość w wysokoczęstotliwościowych sygnałach sterujących, nie wpływa drastycznie na globalny błąd pozycji ze względu na naturalne tłumienie aerodynamiczne i bezwładność płatowca. Ostatecznie stwierdzono, że zarówno zbyt mała, jak i nadmiernie duża czułość autopilota pogarsza jakość sterowania – odpowiednio przez ospałość układu lub jego nadpobudliwość skutkującą agresywną pracę sterów i zbędnym obciążeniem konstrukcji.

\begin{tikzpicture}
  \duck[
    body=red!20!white,
    bill=red!30!white,
    shorthair=red!30!white,
    bunny=red!30!white,
    inear=red!30!white
  ]
  \path (0.1,0.1) rectangle (2.1,2.12);
  \begin{pgfinterruptboundingbox}
    \node[red!20!white,rotate=25,scale=3] at (1.7,1.51) {\textsf{s}};
  \end{pgfinterruptboundingbox}
\end{tikzpicture}
